\section{实验}
\label{sec:experiments}

\subsection{实验协议}

两个主配置都使用 $\lambda=0.75$、stride-32 Key sample、八个 prompt-tail Query、
240-bit plain-scale index、request-local qMSE allocation 和
\cref{eq:budget-intro}。参考配置物化全部 proxy score 并执行 top-$k$；优化配置用
\cref{eq:frozen-quantile-samples}中的 capped sampled-quantile scan 替代它，
request-local 变换、allocation、精确 KV consumer 和预算不变。我们注册
Llama-3.1-8B-Instruct、Qwen3-4B-Instruct 与
Mistral-7B-Instruct-v0.3，完整英文 LongBench、4K--128K RULER，以及 paired
held-out-position PPL。

同步 batch-one 计时分别报告完整 attention、整模型稳态 forward、直接测量的
decode 区间，以及包含 prefill 与初始化的直接请求时间。每个辅助阶段也通过
独立调用直接计时以定位瓶颈，但完整路径绝不由这些阶段数字相加得到。
RTX 3090 数字是已审计诊断；仍需补充 matched H100/A100 测量。

本文主张的范围是：对已经构建且可以复用的 KV cache 加速 decode-time
attention，而不是 sparse prefill。我们将 cold single-use、同一 QK 坐标下的
shared-prefix branching 和 append-only 作为互不混写的计时协议。下文给出已经
审计的 32/64K 切片，其余 H100/128K 网格预注册在
\cref{app:persistent-kv-protocol}；cold request latency 仍需如实报告为局限，不能与
warm-path 加速混成同一结论。

完整 LongBench 表使用 request-local 参考配置。随后的小规模严格筛查直接运行
CUDA 测速所对应的 sampled-quantile 路径。

\begin{table}[t]
\caption{完整 reference-profile LongBench 结果。优化部署配置单独评测。}
\label{tab:main-longbench}
\centering
\small
\begin{tabular}{lrrrr}
\toprule
模型 & Full & \method & 保持率 & Paired 95\% CI \\
\midrule
Llama-3.1-8B & .459398 & .458852 & 99.881\% &
[99.424\%, 100.347\%] \\
\bottomrule
\end{tabular}
\end{table}

该结果覆盖 3,750 个严格配对样本和全部 16 个任务；逐任务分数及 97.615\%
的最低保持率见\cref{app:longbench}。

在 160 个严格同样本 pair 上，早期 Fast/Robust 长度门控路径为 .453676，Full 为
.453019，即保持 100.145\%；paired-sample bootstrap 区间为
[97.915\%,102.051\%]。这在筛查规模上闭合了 reference/deployment 代码路径
差异，但区间仍宽，且 Qasper/MultiNews 分任务保持率分别为 84.28/92.85\%，
因此不能宣称部署质量等价或显著优越。完整 reference 结果仍是主要任务质量证据。

另一个 targeted RULER 审计覆盖 8/16/32/64K 的 13 个任务，每任务、每长度
分别取 2/2/2/1 条。91 个严格 pair 上，Full 与 \method 分别为
.884776/.884135，即保持 99.928\%，paired-bootstrap 区间为
[99.559\%, 100.294\%]。该小实验用于筛查任务族与长度失效模式，不作为
正式大规模 RULER 表；quality harness 中的时间也不进入系统速度结论。

\subsection{等 rate 索引质量}

在 12 个 held-out 32K window 上，使用 240 index bit 和 1,280 个精确 token
时，QK balancing 将 Full-token top-1 agreement 提高 1.95 个百分点，并以
.81\% 延迟代价将 KL 降低 60\%（\cref{tab:key-only-causal}）。Key-MSE 的
下游质量与 logit-MSE 匹配，同时快 2.46\%
（\cref{tab:allocation-causal,tab:allocation-deploy}）。

在 1\% active KV 下，automatic QK-MSE 达到 70.48\% exact-set recall 和
96.36\% oracle-mass recall；其 matched-rate 结果接近 fixed 4-4-2-1。

在 1/2/4\% active KV 下，相对 FIER 的 paired attention-mass 增益分别为
7.08/5.68/4.35 个百分点，区间均排除零 \citep{wang2025fier}。在 59,200 个
condition 上，matched-rate RaBitQ recall/mass 从 53.88/71.32\% 提升到
71.82/76.79\%（\cref{app:mechanism-results}）。

在严格的 160 样本 LongBench 筛查集上，Full macro 为 .45302。相同 active
token schedule 下，门控消融、完整 top-$k$ 的 \method reference、审计过的
FIER RTN-1 g32 reference、Quest-P16、RaBitQ 公式参考、SparQ-R32 selector-only、
完整 SparQ 公式和按公开公式实现的 UNIQUE-P8 保持率分别为
100.15/99.79/100.01/92.22/100.01/99.62/99.95/93.19\%。门控消融与 FIER 的
绝对 macro 差为 .00061，该小筛查集不支持显著优越性结论。这些行只比较
selector 公式与下游质量，不代表官方 kernel 速度。同一 160 pair 上，使用
BinaryPC 发布的 Llama projection 时保持 99.25\%。该门控消融不使用任务 router
或 Full fallback，但每任务 10 条仍不足以宣称显著优越。

\subsection{长度扩展}

\begin{table}[t]
\caption{单张 RTX 3090 上冻结的原生 MHA attention 路径。Full 是不复制 KV
的原生 FP16 SDPA；Fast 包含 selector 与精确稀疏 attention；Robust 还包含
rank-16 INT4 Value-tail 扫描和 consumer。所有行均不含索引构建。}
\label{tab:mha-main}
\centering
\small
\begin{tabular}{rrrrrr}
\toprule
历史长度 & Full & Fast & Robust & Fast 加速 & Robust 加速 \\
\midrule
8K   & .1691 & .1327 & .1560 & 1.27$\times$ & 1.08$\times$ \\
16K  & .3173 & .1900 & .2250 & 1.67$\times$ & 1.41$\times$ \\
32K  & .6163 & .2324 & .2945 & 2.65$\times$ & 2.09$\times$ \\
64K  & 1.2027 & .2646 & .3576 & 4.55$\times$ & 3.36$\times$ \\
128K & 2.3708 & .3722 & .5757 & 6.37$\times$ & 4.12$\times$ \\
\bottomrule
\end{tabular}
\end{table}

冻结的计时形状为 batch one、32 个 Query head、32 个 KV head、$d=128$，FP16
K/V 常驻 GPU。CUDA event 测量包含 10 次 warmup 和 40 次正式迭代；每行取三个
独立 seed 的中位数。稀疏路径包含融合 Query 投影/量化、混合位宽扫描、
sampled-threshold 候选压缩和精确 QK--softmax--AV。15 组 seed--长度对照中，
Fast 与 Robust 的阈值、候选数和候选集合完全相同，唯一差别是是否估计遗漏
Value。kernel 实际分配的 Key-only 与 Key+Value 辅助张量约占 FP16 K+V 的
5.47\% 与 7.09\%；包含摊销 metadata 后的方法级硬上限为 5.86\% 与 7.47\%。

\begin{table}[t]
\caption{与审计后的 packed FIER RTN-1 g32 实现做受控比较。FIER 与 Fast
使用相同 active-token 预算和精确 sparse-attention consumer；Robust 还额外
估计遗漏 Value tail。比值大于一表示 QKSieve 更快。}
\label{tab:mha-fier-samepath}
\centering
\small
\begin{tabular}{rrrrrr}
\toprule
历史长度 & FIER & Fast & Robust & Fast/FIER & Robust/FIER \\
\midrule
8K   & .1816 & .1327 & .1560 & 1.37$\times$ & 1.16$\times$ \\
16K  & .2951 & .1900 & .2250 & 1.55$\times$ & 1.31$\times$ \\
32K  & .4647 & .2324 & .2945 & 2.00$\times$ & 1.58$\times$ \\
64K  & .7329 & .2646 & .3576 & 2.77$\times$ & 2.05$\times$ \\
128K & 1.3322 & .3722 & .5757 & 3.58$\times$ & 2.31$\times$ \\
\bottomrule
\end{tabular}
\end{table}

\Cref{tab:mha-fier-samepath}中，FIER 与 Fast 的模型张量和最终精确 K/V
consumer 都不变，因此主要隔离 packed selector 成本。Robust 还支付 Value-tail
估计成本，是更严格的系统比较。该表支持受控实现对比，不外推到未经实测的
厂商优化 FIER 系统。

\Cref{tab:mha-main}是 attention 子系统结果，不是端到端 decode 主张。它不含
prefill 和历史索引构建，因此我们又在真实 MHA 模型中直接测量相同 Fast 与
Robust 路径。

\begin{table}[t]
\caption{Yarn-Llama-2-7B-128K 上、精确 K/V 常驻时的稳态 greedy decode。
每次生成 64 token，报告最后 48 个 token；数值为直接测得的整模型 ms/token。}
\label{tab:mha-decode}
\centering
\small
\begin{tabular}{rrrrrr}
\toprule
历史长度 & Full & Fast & Robust & Fast 加速 & Robust 加速 \\
\midrule
32K  & 84.18 & 51.56 & 63.97 & 1.63$\times$ & 1.32$\times$ \\
64K  & 144.75 & 52.97 & 65.25 & 2.73$\times$ & 2.22$\times$ \\
128K & 268.17 & 55.41 & 67.38 & 4.84$\times$ & 3.98$\times$ \\
\bottomrule
\end{tabular}
\end{table}

\begin{table}[t]
\caption{同一 MHA 模型和 RTX 3090 上直接计时的 persistent-cache 生命周期。
Cold 包含完整 Robust 索引构建和一个 32-token 分支；warm 复用索引；四分支
均摊把一次构建分给 128 个生成 token。比值均为匹配的 Full/QKSieve 墙钟时间。}
\label{tab:persistent-kv-main}
\centering
\resizebox{\linewidth}{!}{
\begin{tabular}{rrrrrr}
\toprule
历史长度 & Cold & Warm & 四分支均摊 & Append-only & 构建 (s) \\
\midrule
32K & .702$\times$ & 1.322$\times$ & 1.082$\times$ & 1.319$\times$ & 1.759 \\
64K & 1.125$\times$ & 2.221$\times$ & 1.785$\times$ & 2.211$\times$ & 1.998 \\
\bottomrule
\end{tabular}}
\end{table}

32K 的 Full/QKSieve warm decode 为 83.714/63.347 ms/token，64K 为
144.910/65.256 ms/token。32/64K 的一次构建分别由 .724/.735 秒 QK factor、
.518/.603 秒 Key 编码和 .517/.658 秒 ValueSketch 组成。

用实测值计算
$I/(D_f-D_s)$，break-even 约为 87/26 个生成 token，与
\cref{tab:persistent-kv-main}直接测得的 32-token cold 结果一致。

生命周期审计在每个长度覆盖全部 32 层、6 个状态快照和 5 次 rewind。
Key/Value 索引指针与重建计数始终不变；重放 token 和 hash 完全一致；decode 后
两类索引都严格落后精确 cache 一个注册 token。该实验只在分支与连续生成之间
复用同一个 request-local QK transform。它没有证明新 prompt 改变 Query 统计后
仍能复用旧索引；这种情况必须重新构建 request-local 索引。

原生 128K 的修正同 seed 六主题审计推翻了早期结论。QKSieve-Fast 的 pooled
PPL 保持率仅为 96.319\%（case macro 96.373\%），政治主题最差为 91.100\%。
一个医学归因窗口中，Exact-QK top-1,280 也只有 94.74\%，预算加倍后为
96.24\%；因此问题不只是低比特 Key 错排。额外保留 128/256 个 recent token
使弱项 macro 变化不到 .1 个点。

QKSieve-Robust rank-16 INT4 的 $\alpha=.5$ 在六个同 seed case 上冻结，并在
六个新 seed 上独立评估。12 个 case 的 pooled PPL 保持率为 99.805\%，
case-bootstrap 95\% 区间为 [99.100\%,100.499\%]，case macro 为 99.812\%，
最差为 96.980\%。逐 token 辅助 rate 总计 7.471\%。较早的 GQA-4 子系统和
整模型速度保留在\cref{app:system-diagnostics}中，不与\cref{tab:mha-main}的
原生 MHA 速度合并。因此 always-on Robust 是冻结主方法，Fast 只作为去掉
ValueSketch 的因果消融。

长度门控消融和 256K 坐标外推失效见\cref{app:system-diagnostics}；它们不改变
冻结的 always-on Robust 合同及其注册的 4K--128K 适用范围。
